\chapter{Babies and Toddlers}

\section*{Definition and Overview}
The care of babies and toddlers covers the health and development of children in the first few years of life. Babies are usually taken to mean infants under one year, and toddlers are children from around one to three years. This is a time of rapid growth, of major developmental milestones, and of frequent illnesses, as young children build their immune systems and explore the world around them.

Children are not simply small adults: their bodies and immune systems are immature, they cannot always describe how they feel, and their health depends heavily on their parents and carers. Most illnesses in this age group are minor and self-limiting, but a small number are serious. Recognising the child who is seriously unwell is the central task of parents and health professionals alike.

\section*{Epidemiology}
Most babies are healthy, but the first year of life carries the highest risk of serious illness of any childhood period. The main causes of death in infancy are conditions present at birth, prematurity and sudden infant death syndrome (SIDS). Beyond the first year, accidental injury becomes a leading cause of death in developed countries.

Common reasons babies and toddlers see a doctor include viral respiratory infections, ear infections, gastroenteritis and fevers. Bronchiolitis is a common cause of hospital admission in winter. Immunisation has dramatically reduced diseases such as measles, whooping cough and meningitis in vaccinated populations, though the age group remains most vulnerable to these infections when outbreaks occur.

\section*{Aetiology and Pathophysiology}
Illness in this age group has many causes, and the same symptom often has different causes at different ages.

\begin{itemize}
    \item \textbf{Infections:} viruses causing colds, bronchiolitis, croup and gastroenteritis; bacteria causing ear infections, urinary infections and pneumonia
    \item \textbf{Congenital conditions:} heart defects, metabolic problems and other conditions present from birth
    \item \textbf{Developmental:} teething, colic and normal behavioural stages such as tantrums
    \item \textbf{Allergic conditions:} eczema, food allergy and cow's milk protein allergy
    \item \textbf{Nutritional:} iron deficiency, vitamin D deficiency and under- or over-nutrition
    \item \textbf{Injury:} falls, burns, choking, poisoning and drowning
\end{itemize}

Babies are especially vulnerable to infection because their immune systems are still learning to fight germs, and they cannot regulate their body temperature or fluid balance as efficiently as older children. Their airways are small, so swelling from even a mild respiratory infection can make breathing difficult.

\section*{Clinical Features}
Parents should look for the general signs that a baby or toddler is unwell, as well as specific symptoms.

\begin{itemize}
    \item Fever, which is the most common reason for concern
    \item Reduced feeding or drinking, or a nappy that stays dry
    \item Lethargy, sleepiness or unusual irritability
    \item Breathing fast, noisily or with the ribs pulling in (recession)
    \item Vomiting, diarrhoea or tummy pain
    \item A rash, particularly one that does not fade when pressed
    \item A fit (seizure) or unusual movements
    \item Floppiness, or a baby who is difficult to wake
    \item Slow weight gain or missed developmental milestones
\end{itemize}

\section*{Differential Diagnosis}
Many symptoms in this age group have more than one possible cause, and distinguishing them matters because some are dangerous.

\begin{itemize}
    \item Viral illness versus serious bacterial infection, such as meningitis or a urine infection
    \item Bronchiolitis versus pneumonia versus asthma or wheeze
    \item Croup versus choking versus an inhaled object
    \item Gastroenteritis versus other causes of vomiting, such as a bowel obstruction
    \item Colic or reflux versus organic disease in a young infant
    \item A normal developmental variation versus a condition needing early intervention
    \item Accidental injury versus injury that may have been inflicted
\end{itemize}

\section*{When to Seek Medical Care}
See a doctor for any baby who has a fever under three months old, is feeding much less than usual, is passing very little urine, has a symptom that is not improving after 48 hours, or has any concern about breathing, rash, growth or development. When in doubt, trust your instincts and seek advice.

\textbf{Call 000 (or local emergency services) for a baby who looks seriously ill:}
\begin{itemize}
    \item Is difficult to wake, floppy or not responding
    \item Has difficulty breathing, grunting or pauses in breathing
    \item Has a fit that does not stop, or blue lips or skin
    \item Has a rash that does not fade when pressed
    \item Has blood in vomit or stools, or a sunken fontanelle with dry nappies
\end{itemize}

\section*{Diagnosis and Investigations}
Assessment starts with the story from the parent or carer, then a careful examination, with tests used only when needed.

\begin{itemize}
    \item \textbf{History:} onset and course of symptoms, feeding, wet nappies, activity and contact with unwell people
    \item \textbf{Examination:} temperature, breathing and heart rate, colour, alertness, and the skin, ears, throat and chest
    \item \textbf{Growth measurement:} weight, length and head circumference plotted on growth charts
    \item \textbf{Developmental assessment:} milestones for age, such as smiling, sitting and walking
    \item \textbf{Investigations:} urine dipstick and culture, throat swabs or blood tests when a serious infection is suspected
    \item \textbf{Imaging:} a chest X-ray or other scans only when clinically indicated
    \item \textbf{Referral:} to a paediatrician for complex or ongoing problems
\end{itemize}

\section*{Management}
Most care is supportive and happens at home, with hospital care reserved for serious illness.

\begin{itemize}
    \item \textbf{Supportive care at home:} rest, small frequent feeds and fluids to maintain hydration
    \item \textbf{Paracetamol or ibuprofen:} for fever or discomfort, at the weight-appropriate dose and under guidance
    \item \textbf{Oral rehydration solution:} for gastroenteritis, given in small frequent amounts
    \item \textbf{Specific treatment:} antibiotics for bacterial infection, or inhalers and other medicines for conditions such as asthma
    \item \textbf{Hospital care:} for dehydration, breathing difficulty or suspected serious infection, often with fluids and oxygen
    \item \textbf{Family-centred care:} clear advice on when to return, since babies can deteriorate quickly
\end{itemize}

\section*{Complications}
\begin{itemize}
    \item Dehydration and electrolyte disturbance, especially with gastroenteritis
    \item Febrile seizures with high fevers
    \item Serious bacterial infections such as meningitis or sepsis, which can be rapidly progressive
    \item Breathing failure with severe bronchiolitis or croup
    \item Long-term effects of iron deficiency on development
    \item Injuries and their consequences, from falls, choking, poisoning and drowning
    \item Missed developmental delay if not identified early
\end{itemize}

\section*{Prognosis and Outcomes}
Most babies and toddlers recover quickly and completely from common illnesses, and the vast majority grow and develop normally. The outlook is excellent with timely care: serious infections are treatable when caught early, and injuries are largely preventable. The most dangerous period is the first year, so health checks and safety measures in infancy matter most. Children with developmental delay do far better when it is recognised early and support is started promptly.

\section*{Prevention and Self-Care}
\begin{itemize}
    \item Follow the immunisation schedule, starting at birth
    \item Put babies to sleep on their back, alone in a cot with no loose bedding, to reduce the risk of SIDS
    \item Breastfeed if possible, and support healthy feeding and weaning
    \item Use an approved, correctly fitted car seat on every trip
    \item Make the home safe: stair gates, window locks, secured furniture, and locked medicines and poisons
    \item Keep small objects, hot drinks and pools away from young children, and supervise them constantly near water
    \item Attend scheduled health checks and growth and development reviews
    \item Follow safe sleep, safe food and safe medicine practices, and ask for advice whenever uncertain
\end{itemize}

\section*{Sources and Further Reading}
The following reputable organisations provide reliable, current information on the health of babies and toddlers.

\begin{itemize}
    \item Raising Children Network (Australia). \textit{Evidence-based information for parents on babies and toddlers.} https://raisingchildren.net.au
    \item National Health Service. \textit{Health information for parents of babies and young children.} https://www.nhs.uk/conditions/baby/
    \item Centers for Disease Control and Prevention. \textit{Information on infant and toddler development and safety.} https://www.cdc.gov/ncbddd/childdevelopment/
    \item Australian Government Department of Health and Aged Care. \textit{Information on maternal, child and family health services.} https://www.health.gov.au
    \item Red Nose Australia. \textit{Information on safe sleep and reducing the risk of SIDS.} https://rednose.org.au
    \item World Health Organization. \textit{Information on child health, growth and development.} https://www.who.int/health-topics/child-health
\end{itemize}
